\section*{Capítulo \ref{ch:b1:earthquake-source} --- La fuente del terremoto}

\begin{solution}{exo:b1:earthquake-source:1}
(a) $\log_{10}(\num{4e17}) = 17.60$, así que $M_w = \tfrac23(17.60 - 9.1) = 5.67$.
(b) $\log_{10}(\num{5.59e22}) = 22.75$ y $M_w = \tfrac23(22.75 - 9.1) =
9.10$. (c) $M_0 = 10^{1.5\times7.0 + 9.1} = 10^{19.6} = \qty{3.98e19}{N.m}$.
\end{solution}

\begin{solution}{exo:b1:earthquake-source:2}
(a) $T = D/\dot s = \qty{3}{m}/\qty{0.025}{m.yr^{-1}} = \qty{120}{yr}$.
(b) Con $c = 0.5$, $T = D/(c\dot s) = \qty{240}{yr}$: la reptación absorbe la mitad
del movimiento, y los terremotos, del mismo deslizamiento, llegan la mitad de a menudo.
\end{solution}

\begin{solution}{exo:b1:earthquake-source:3}
Por \cref{prop:b1:earthquake-source:moment-steps}, el cociente de los momentos es
$10^{1.5\times(7.4 - 6.2)} = 10^{1.8} = 63.1$: hacen falta 63 terremotos de
magnitud 6.2 para liberar el momento de uno de magnitud 7.4.
\end{solution}

\begin{solution}{exo:b1:earthquake-source:4}
Con $b = 1$, cada unidad de magnitud divide el recuento por 10: 10 terremotos de
magnitud 5 o más al año, 1 de magnitud 6 o más, y 0.1 de magnitud 7 o más, es decir, uno
cada 10 años de media.
\end{solution}

\begin{solution}{exo:b1:earthquake-source:5}
(a) $A = \num{4.3e5} \times \qty{1.5e4}{m} = \qty{6.45e9}{m^2}$, y
$M_0 = \num{3e10}\times\num{6.45e9}\times4 = \qty{7.74e20}{N.m}$;
$\log_{10}M_0 = 20.89$ y $M_w = \tfrac23(20.89 - 9.1) = 7.86$. (b) Una profundidad de
\qty{10}{km} divide el área y el momento por 1.5, lo que baja la magnitud en
$\tfrac23\log_{10}1.5 = 0.12$, hasta 7.74. (c) La primera, 7.86, redondea al
7.9 del catálogo; la segunda queda un poco baja. La magnitud no es muy sensible a la
profundidad: un factor de 1.5 en el área equivale solo a 0.12 de magnitud.
\end{solution}

\begin{solution}{exo:b1:earthquake-source:6}
(a) $D = M_0/(\mu A) = \num{5.59e22}/(\num{3e10}\times\num{4e5}\times\num{1.5e5})
= \qty{31.1}{m}$. (b) $\qty{31.1}{m}/\qty{0.083}{m.yr^{-1}} = \qty{374}{yr}$.
(c) Con la fracción bloqueada 0.8, $31.1/(0.8\times0.083) = \qty{468}{yr}$; con
0.5, \qty{749}{yr}. Cuanto menos bloqueada está la interfase, más tarda en acumularse
el deslizamiento de un gran terremoto; el deslizamiento medio hallado aquí, con un
rectángulo de \qtyproduct{400 x 150}{km} y $\mu = \qty{30}{GPa}$, es en sí mismo
solo una estimación, pues el deslizamiento real estuvo lejos de ser uniforme.
\end{solution}

\begin{solution}{exo:b1:earthquake-source:7}
(a) El \omterm{def:b1:geological-maps:attitude}{rumbo} es $\phi = 329$, la \omterm{def:b1:geological-maps:attitude}{dirección de buzamiento} $a = 059$, el \omterm{def:b1:geological-maps:attitude}{buzamiento} $\delta =
\ang{8}$. Entonces $\mathbf s = (\sin329, \cos329, 0) = (-0.5150, 0.8572, 0)$,
$\mathbf u = (-\cos8\sin59, -\cos8\cos59, \sin8) = (-0.8488, -0.5100, 0.1392)$,
$\mathbf n = (\sin8\sin59, \sin8\cos59, \cos8) = (0.1193, 0.0717, 0.9903)$,
y con $r = \ang{110}$,
\[
  \mathbf d = \cos110\,\mathbf s + \sin110\,\mathbf u = (-0.6215, -0.7724, 0.1308) .
\]
(b) El \omterm{def:b1:earthquake-source:focal-mechanism}{plano auxiliar} tiene la normal $\mathbf d$ y el deslizamiento $\mathbf n$.
Tomando la normal hacia arriba, buza $\arccos 0.1308 = \ang{82.5}$ hacia
$\arctan(-0.6215/-0.7724) = 218.8$, así que su \omterm{def:b1:geological-maps:attitude}{rumbo} es $128.8$; el deslizamiento
$\mathbf n$ proyectado sobre sus vectores de \omterm{def:b1:geological-maps:attitude}{rumbo} y de máxima pendiente da el \omterm{def:b1:faults-folds-fabrics:fault-elements}{cabeceo}
\ang{87.2}. Redondeado: $129/83/87$, el segundo plano de la solución publicada.
(c) $\mathbf P = (\mathbf n - \mathbf d)/\sqrt2$ tiene tendencia $221$ e \omterm{def:b1:field-lab-methods-1:trend-plunge}{inmersión}
\ang{37}; $\mathbf T = (\mathbf n + \mathbf d)/\sqrt2$ tiene tendencia $036$ e \omterm{def:b1:field-lab-methods-1:trend-plunge}{inmersión}
\ang{52} (publicado: $222/38$ y $36/52$). (d) El plano somero, que buza \ang{8} hacia $059$, a lo largo de la
\omterm{def:b1:active-margins:slab}{losa en subducción}: una \omterm{def:b1:faults-folds-fabrics:fault-types}{falla inversa} con \omterm{def:b1:faults-folds-fabrics:fault-elements}{cabeceo}
de \ang{110}, y, al buzar menos de \ang{45}, un \omterm{def:b1:faults-folds-fabrics:fault-types}{cabalgamiento}, la interfase de placas
de la zona de subducción.
\end{solution}

\begin{solution}{exo:b1:earthquake-source:8}
(a) El plano buza hacia $135$, así que $\mathbf n = (\sin135, \cos135, 0) =
(0.707, -0.707, 0)$; un deslizamiento \omterm{def:b1:faults-folds-fabrics:fault-types}{dextral} tiene $r = \ang{180}$, $\mathbf d = -\mathbf s
= (-0.707, -0.707, 0)$. Hacia el norte, $\mathbf x = (0, 1, 0)$:
$(\mathbf n\cdot\mathbf x)(\mathbf d\cdot\mathbf x) = (-0.707)(-0.707) > 0$, un
empuje. Al este: $(0.707)(-0.707) < 0$, un tirón. Al sur, empuje; al oeste, tirón.
(b) $\mathbf P = (\mathbf n - \mathbf d)/\sqrt2 = (1, 0, 0)$: el \omterm{def:b1:earthquake-source:focal-mechanism}{eje P} es
horizontal y tiene tendencia $090$--$270$; $\mathbf T = (0, -1, 0)$ tiene tendencia
$000$--$180$. (c) El plano vertical de \omterm{def:b1:geological-maps:attitude}{rumbo} $135$--$315$ tiene la normal
$\mathbf d$. Con el deslizamiento $\mathbf n$, el bloque de su lado suroeste, hacia el que apunta
$\mathbf d$, se mueve hacia $135$; visto desde el bloque nordeste, mirando al suroeste, eso
queda a la izquierda: un deslizamiento \omterm{def:b1:faults-folds-fabrics:fault-types}{sinistral}. Por
\cref{prop:b1:earthquake-source:auxiliary}, este deslizamiento da los mismos primeros
movimientos en todas las direcciones.
\end{solution}

\begin{solution}{exo:b1:earthquake-source:9}
(a) $V = \qty{40}{mm.yr^{-1}}$ y $v = V/4$ a \qty{10}{km}, así que por
\cref{prop:b1:earthquake-source:profile}(b) $D = \qty{10}{km}$; directamente,
$D = x/\tan(\pi v/V) = 10/\tan(\pi/4) = \qty{10}{km}$. (b) $v(30) =
(40/\pi)\arctan 3 = \qty{15.9}{mm.yr^{-1}}$. (c) $V/(\pi D) =
\num{0.040}/(\pi\times\num{10000}) = \num{1.27e-6}$ por año.
\end{solution}

\begin{solution}{exo:b1:earthquake-source:10}
(a) Los momentos de los intervalos siguen siendo $C q^k$ con $q = 10^{1.5 - b}$, pero ahora $q < 1$: el momento de los intervalos
decrece hacia las magnitudes grandes, y el intervalo inferior, $k = 0$, lleva la fracción $C/S = (1 - q)/(1 - q^{K+1}) > 1 - q$, la imagen
especular de \cref{prop:b1:earthquake-source:largest}. Los terremotos más pequeños contados dominarían entonces, y el total crecería sin
límite a medida que se incluyeran otros cada vez más pequeños. (b) $1 - 10^{-0.7} = 0.80$: al menos
\qty{80}{\percent}. (c) $\log_{10}N = 8.20 - 1.01\times9 = -0.89$, así que $N = 0.13$
al año y unos 6.4 en cincuenta años, tres veces los dos observados. La recta,
ajustada entre las magnitudes 6 y 7.5, sobrestima los mayores terremotos; o bien la ley
se curva hacia abajo en lo alto, como sugiere el tamaño finito de las fallas, o cincuenta
años son pocos para muestrear sucesos que llegan una vez cada decenas de años. Ambas
cosas son posibles con dos sucesos.
\end{solution}

\begin{solution}{exo:b1:earthquake-source:11}
(a) En 150 años un poste en $x$ se mueve $150\,v(x) = 150\times(34/\pi)
\arctan(x/15)$ milímetros respecto del poste sobre la falla: \qty{0.11}{m} a
\qty{1}{km}, \qty{1.28}{m} a \qty{15}{km}, \qty{2.08}{m} a \qty{50}{km},
hacia el sur para una falla \omterm{def:b1:faults-folds-fabrics:fault-types}{dextral} vista desde el oeste. (b) El déficit de deslizamiento es
$s = \qty{0.034}{m.yr^{-1}}\times\qty{150}{yr} = \qty{5.1}{m}$. Durante el
terremoto los postes se mueven $(s/\pi)\arctan(D/x)$: \qty{2.44}{m} a
\qty{1}{km}, \qty{1.28}{m} a \qty{15}{km}, \qty{0.47}{m} a \qty{50}{km}.
(c) Las sumas son $0.11 + 2.44$, $1.28 + 1.28$ y $2.08 + 0.47$, todas
\qty{2.55}{m} $= s/2$: tras el terremoto la cerca vuelve a estar recta en el lado
este, desplazada $s/2$, y el lado oeste, por simetría, $-s/2$.
\end{solution}

\begin{solution}{exo:b1:earthquake-source:12}
(a) $\mathbf P = (0, 1, 0)$ y $\mathbf T = (1, 0, 0)$. El par
$\mathbf n = (\mathbf T + \mathbf P)/\sqrt2 = (0.707, 0.707, 0)$,
$\mathbf d = (\mathbf T - \mathbf P)/\sqrt2 = (0.707, -0.707, 0)$ cumple
$\mathbf P = (\mathbf n - \mathbf d)/\sqrt2$ y $\mathbf T = (\mathbf n +
\mathbf d)/\sqrt2$. Ambos vectores son horizontales, así que ambos \omterm{def:b1:earthquake-source:focal-mechanism}{planos nodales} son
verticales: el plano de normal $\mathbf n$ (hacia $045$) tiene \omterm{def:b1:geological-maps:attitude}{rumbo} $135$--$315$,
el plano de normal $\mathbf d$ (hacia $135$) tiene \omterm{def:b1:geological-maps:attitude}{rumbo} $045$--$225$. (b) En el
plano de \omterm{def:b1:geological-maps:attitude}{rumbo} $315$, el bloque del lado nordeste se mueve según $\mathbf d$,
hacia $135$: visto desde el otro bloque, a la derecha, un deslizamiento \omterm{def:b1:faults-folds-fabrics:fault-types}{dextral}. En el
plano de \omterm{def:b1:geological-maps:attitude}{rumbo} $045$, el bloque del lado sudeste se mueve según $\mathbf n$,
hacia $045$: un deslizamiento \omterm{def:b1:faults-folds-fabrics:fault-types}{sinistral}. (c) Una falla nueva forma $\ang{45} - \phi/2 =
\ang{30}$ con $\sigma_1$: con $\sigma_1$ con tendencia $000$, las fallas
tendrían \omterm{def:b1:geological-maps:attitude}{rumbo} $030$ (\omterm{def:b1:faults-folds-fabrics:fault-types}{sinistral}) y $330$ (\omterm{def:b1:faults-folds-fabrics:fault-types}{dextral}), a \ang{15} de los \omterm{def:b1:earthquake-source:focal-mechanism}{planos nodales}.
El \omterm{def:b1:earthquake-source:focal-mechanism}{eje P} queda fijado a \ang{45} de los \omterm{def:b1:earthquake-source:focal-mechanism}{planos nodales} por la geometría, sea cual sea
la fricción de la roca, así que coincide con $\sigma_1$ solo si la falla deslizó
a \ang{45} de él. (d) La falla de \omterm{def:b1:geological-maps:attitude}{rumbo} $040$ está a \ang{5} del \omterm{def:b1:earthquake-source:focal-mechanism}{plano nodal} de \omterm{def:b1:geological-maps:attitude}{rumbo} $045$:
ese plano es probablemente la falla, y el deslizamiento sobre él es
\omterm{def:b1:faults-folds-fabrics:fault-types}{sinistral}.
\end{solution}

\begin{solution}{pb:b1:earthquake-source:1}
\textbf{1.} Los intervalos son $1290 - 1050 = 240$, $1480 - 1290 = 190$ y
$1712 - 1480 = 232$ años; su media es \qty{220.7}{yr}.

\textbf{2.} El desplazamiento medio es $(7.0 + 6.0 + 5.5 + 7.5)/4 = \qty{6.5}{m}$, y
$T = 6.5/0.030 = \qty{216.7}{yr}$, cerca del intervalo medio.

\textbf{3.} Los desplazamientos de los tres últimos terremotos, $6.0 + 5.5 + 7.5 =
\qty{19.0}{m}$, se acumularon entre 1050 y 1712, en 662 años: $\qty{19.0}{m}
/\qty{662}{yr} = \qty{28.7}{mm.yr^{-1}}$, un \qty{4}{\percent} menos que la
velocidad geodésica de \qty{30}{mm.yr^{-1}}. Ambas concuerdan dentro de la incertidumbre de las
fechas y de los desplazamientos: el segmento libera casi toda su carga en sus grandes terremotos.

\textbf{4.} $2025 - 1712 = 313$ años, y el déficit es
$\qty{0.030}{m.yr^{-1}}\times\qty{313}{yr} = \qty{9.39}{m}$.

\textbf{5.} El déficit supera cada desplazamiento de la trinchera, el mayor de los cuales es
\qty{7.5}{m}, y el tiempo desde 1712 supera el intervalo más largo, 240 años.
En ese sentido el segmento va retrasado; no significa que el terremoto sea inminente, solo que la falla
ya contiene más deslizamiento del que ha liberado nunca de una vez en el registro.

\textbf{6.} Los intervalos van de 190 a 240 años y los desplazamientos de
\qtyrange{5.5}{7.5}{m}: el ciclo no es periódico, y la resistencia de la falla en el
próximo terremoto es desconocida. La trinchera da una media y una dispersión, a partir
de las cuales puede estimarse una probabilidad, no una fecha.

\textbf{7.} $A = \qty{250}{km}\times\qty{15}{km} = \qty{3750}{km^2} =
\qty{3.75e9}{m^2}$.

\textbf{8.} $M_0 = \num{3e10}\times\num{3.75e9}\times\num{9.39} =
\qty{1.06e21}{N.m}$.

\textbf{9.} $M_0 = \num{3e10}\times\num{3.75e9}\times\num{7.5} =
\qty{8.44e20}{N.m}$, así que $M_w = \tfrac23(20.926 - 9.1) = 7.88$.

\textbf{10.} El cociente es $9.39/7.5 = 1.25$, que equivale a $\tfrac23\log_{10}1.25 =
0.07$ de magnitud.

\textbf{11.} La mitad del área reduce a la mitad el momento y baja la magnitud en
$\tfrac23\log_{10}2 = 0.20$: $M_w = 7.75$.

\textbf{12.} $\mu A V = \num{3e10}\times\num{3.75e9}\times\num{0.030} =
\qty{3.38e18}{N.m}$ por año.

\textbf{13.} $b = 1$. $N(5) = 10^{4.3 - 5} = 0.20$ al año y $N(6) = 0.020$ al
año.

\textbf{14.} $N(7) = 0.0020$ al año: uno cada 500 años; $N(8) = 0.00020$ al
año: uno cada \num{5000} años.

\textbf{15.} La ley da un terremoto de magnitud 7.95 o más cada
$1/10^{4.3 - 7.95} = \num{4500}$ años, y la trinchera uno cada 220 años, veinte
veces más a menudo. La ley se ajusta a los muchos terremotos pequeños de todas las
fallas de la región; extrapolada a los grandes terremotos de un solo segmento, los
subestima mucho. La recurrencia de los grandes terremotos de una falla se toma de su
propio registro y de su carga, no de los pequeños.

\textbf{16.} $\num{5.5e17}/\num{3.38e18} = 0.16$: los terremotos moderados liberan
\qty{16}{\percent} del momento que acumula el segmento.

\textbf{17.} Por \cref{prop:b1:earthquake-source:moment-steps},
$10^{1.5\times(7.95 - 6)} = 10^{2.92} \approx 840$ veces más.

\textbf{18.} Con $b = 1$, el momento de los terremotos lo domina el mayor: lo que hagan
los moderados, liberan una parte menor, aquí \qty{16}{\percent}, y el resto,
\qty{84}{\percent}, el déficit de deslizamiento, espera al gran terremoto.

\textbf{19.} $v(-x) = -v(x)$ para cada par, así que el perfil es antisimétrico.
Lejos de la falla las velocidades tienden a $\pm V/2 = \pm\qty{15}{mm.yr^{-1}}$.

\textbf{20.} $D = x/\tan(\pi v/V) = 15/\tan(\pi\times7.5/30) = 15/\tan\ang{45} =
\qty{15}{km}$: a la \omterm{def:b1:earthquake-source:locking-depth}{profundidad de bloqueo} la velocidad es $V/4$.

\textbf{21.} A \qty{5}{km}: $D = 5/\tan(\pi\times3.1/30) = 5/\tan\ang{18.6} =
\qty{14.9}{km}$; a \qty{40}{km}: $D = 40/\tan\ang{69.6} = \qty{14.9}{km}$. Las
velocidades se dan con \qty{0.1}{mm.yr^{-1}}, y solo ese redondeo mueve $D$ unos
\qty{0.1}{km}; con datos reales, el ruido de cada estación sería mayor, y la
\omterm{def:b1:earthquake-source:locking-depth}{profundidad de bloqueo} se ajustaría a todas las estaciones a la vez.

\textbf{22.} La \omterm{def:b1:earthquake-source:locking-depth}{profundidad de bloqueo}, \qty{15}{km}, iguala la profundidad de los
terremotos más profundos. La parte bloqueada de la falla es su parte frágil, la
\omterm{def:b1:stress-strain-rheology:strength-envelope}{capa sismogénica}, donde el rozamiento la mantiene entre terremotos y falla en ellos; por debajo, la roca más
caliente fluye lentamente (\cref{ch:b1:stress-strain-rheology}) y la falla desliza de manera constante, que es el deslizamiento a la
velocidad $V$ por debajo de $D$ del modelo.

\textbf{23.} $V/(\pi D) = \num{0.030}/(\pi\times\num{15000}) = \num{6.4e-7}$
por año.

\textbf{24.} Con $M_0 = \qty{1.056e21}{N.m}$, $\log_{10}M_0 = 21.024$ y
$M_w = \tfrac23(21.024 - 9.1) = 7.949$. Con una precisión de 0.01, el terremoto que liberaría
todo el déficit del segmento tiene la \omterm{def:b1:earthquake-source:mw}{magnitud de momento}
\textbf{7.95}.
\end{solution}
